% W6 draft -- workshop paper. Swap to the ICLR 2027 workshop style file when
% released; plain article until then. Every number cites its RESEARCH_LOG
% entry in a comment (Lnnn), so a reviewer's question traces to a logged run.
% \pending{...} marks what S1 (PREREG 2026-09-29) must fill; do not submit
% while any remains:  grep -n pending paper/main.tex
\documentclass[11pt]{article}
\usepackage[margin=1in]{geometry}
\usepackage{amsmath,amssymb,booktabs,graphicx,xcolor,hyperref}
\newcommand{\pending}[1]{\textcolor{red}{[#1]}}

\title{Neither Memory nor Continuous Time: A Pre-Registered Negative Result\\
for Learned Adaptive-Therapy Policies under Burden-Only Observation}
\author{Anonymous}
\date{}

\begin{document}
\maketitle

\begin{abstract}
Continuous-time recurrent policies are often proposed for control problems
whose dynamics are themselves differential equations, and irregularly observed
ones in particular. We test this inductive-bias claim on adaptive cancer
therapy, where a policy doses a virtual patient from noisy total tumour burden
alone. One RL algorithm (PQN) trains a memoryless MLP, a parameter-matched
LSTM, an LSTM with the inter-visit gap as input, and an ODE-RNN on a
Lotka--Volterra competition model. We compare them on paired virtual patients,
with confidence intervals that resample both training seeds and patients. Every
comparison was pre-registered. The ODE-RNN shows no edge over the best
discrete arm in range, under a parameter shift, or with up to 50\% of visits
missed. At regular visits the memoryless MLP learns a burden-threshold switch
and only ties a one-line rule. When half of the visits are missed, no learned
policy beats dosing at the maximum tolerated dose throughout. \pending{S1:
whether both hold on Strobl et al.'s published parameter ranges.} We also report
four ways this evaluation misled us before it was valid, each caught by a
pre-registered check.
\end{abstract}

\section{Introduction}
Adaptive therapy doses a tumour to keep a drug-sensitive population alive as a
competitor of the resistant one, rather than to eradicate it. Reinforcement
learning has been used to schedule it: on a Lotka--Volterra (LV) model with
burden-only observation, A3C learns a treat-above-threshold rule close to the
clinical algorithm \cite{gallagher2024}. Recurrence has been reported to help
chemotherapy control when part of the state is hidden \cite{kiram2026}. A
natural next claim is that a policy whose hidden state is itself an ODE
\cite{rubanova2019} should carry the right inductive bias. It should be most
useful when the biology shifts or when visits are irregular.

We test that claim as a falsifiable, pre-registered comparison, and it fails.
Our contributions are:
\begin{itemize}
\item A strict ablation: one trainer, one environment and one evaluation, with
only the Q-network architecture swapped. The recurrent arms are matched at
${\sim}3{,}100$ Q-path parameters.
\item Three negative results with seed+patient confidence intervals: no
continuous-time edge in range (Tier~0), under a parameter shift (Tier~1), or
under missed visits (Tier~4).
\item One finding about the task: as visits go missing, the value of feedback
collapses below that of a fixed maximum-dose schedule.
\item Four evaluation traps that each produced a publishable-looking wrong
answer on the way (\S\ref{sec:pitfalls}).
\end{itemize}

\section{Setup}
\paragraph{Virtual patients.} Sensitive $S$ and resistant $R$ cells compete for
a shared capacity $K$:
\begin{align*}
\dot S &= r_S S\,(1 - \tfrac{S+R}{K})(1 - d_D u) - d_S S, &
\dot R &= r_R R\,(1 - \tfrac{S+R}{K}) - d_R R,
\end{align*}
with drug effect $u\in[0,1]$ from one-compartment PK. Patients are
Latin-hypercube samples of $\theta$ and of the initial resistant fraction $f_R$
and fill level $n_0$. The primary regime follows Strobl et al.\ \cite{strobl2021}
(Table~1): $r_S = 0.027\,\mathrm{d}^{-1}$, $r_R/r_S \in [0.5, 1]$, one turnover
$d_T/r_S\in[0,0.5]$, $d_D=1.5$, $n_0\in[0.1,0.75]$, $f_R\in[10^{-3},10^{-1}]$.
A second, provisional regime has a much stronger cost of resistance (median
$r_R/r_S = 0.30$). It is reported as a robustness check, not as a parameterisation
from the literature. % L703
\paragraph{Episodes.} 52 weekly visits (357\,d). The policy sees only
$\log(\text{burden}/\text{baseline})$ with 5\% log-normal noise, never $S$ or
$R$, and picks a dose in $\{0, \tfrac12, 1\}$. Progression (TTP) is a 20\% rise
over the running nadir, confirmed over 28\,d. The reward is
survival-per-day minus a burden penalty minus $w\cdot$dose; we sweep
$w\in\{0.05,0.1,0.2,0.4\}$.
\paragraph{Tiers.} Tier~0 uses the training ranges. Tier~1 scales
$r_S, r_R$ by 1.5 and $K$ by 0.6, outside the training box. Tier~4 misses each
visit with probability $p$ i.i.d.; a missed visit gives no observation and
no decision, and the last dose continues (``hold'').
\paragraph{Arms.} MLP (width 441), LSTM (26 units), LSTM+$\Delta t$ (the gap as
a second input, the standard discrete baseline for irregular series), and
ODE-RNN \cite{rubanova2019}. The ODE-RNN is a GRU jump at each observation,
with a tanh-bounded latent flow between them in fixed $7/8$\,d steps, so gaps of
whole weeks compose exactly. All use a LayerNorm Q-readout, trained by PQN
\cite{gallici2024} with its published RNN settings, untuned here.
\paragraph{Statistic.} Per patient,
$U=\min(\mathrm{TTP},357)-w\cdot\text{exposure}$, the trained objective's two
main terms. $\Delta(A,B)$ is the mean over $w$ (and, at Tier~4, over
$p\in\{0.2,0.5\}$) of the paired mean $U_A-U_B$. The 95\% CI comes from 2000
bootstrap resamples of \emph{training seeds and patients}. We report capped
median TTP with IQR, median exposure, and the burden integral (a metastasis
proxy) alongside it. Rules are the references: no treatment, MTD (full dose
throughout), AT50 (the clinical adaptive-therapy trial rule), and thr0.8 (full
dose above $0.8\times$ baseline).

\section{Results}
\pending{S1 Part A table: ODE$-$MLP, ODE$-$LSTM, LSTM$-$MLP at Tiers 0/1, Strobl
regime, with the provisional regime (W3c, L565/L596) beside it.}

\paragraph{No continuous-time edge at regular visits.} On the provisional
regime, LSTM$-$MLP is $-82.9$\,d $[-106.8,-58.7]$ and ODE$-$LSTM is $-15.7$\,d
$[-43.0,+21.8]$ at Tier~1 (L565, L596). Tier~0 lies within 10\,d of Tier~1, so
the gap is not a robustness effect. The MLP's dose is a function of the current
observation in $\geq 0.94$ of visits: dose off below ${\sim}0.8\times$ baseline,
full dose above (L461, L554). It ties thr0.8 (246.8 vs 245.9\,d mean $U$).
\pending{Strobl regime.}

\paragraph{Missed visits.} \pending{S1 Part B.} On the provisional regime
(W4b, L668), with no cell excluded, ODE$-$MLP is $-6.1$ $[-34.7,+12.9]$ and
ODE$-$MTD is $-23.7$ $[-47.5,-5.8]$. The ODE-RNN degrades least of the learned
arms ($U$ 246$\to$229 from $p=0$ to $0.5$; MLP 281$\to$217), but at $p=0.5$
every learned arm is below MTD (best $-29.1$ $[-55.8,-8.2]$). The mechanism
under hold is that a threshold policy that switched off keeps the drug off
through the gap while the tumour regrows.

\section{Four ways the comparison misled us}\label{sec:pitfalls}
\begin{enumerate}
\item \textbf{Open-loop collapse.} Under PPO, the ODE-RNN gave all 200
patients one identical exposure in 4 of 6 cells. The first ``architecture''
result compared feedback with a fixed schedule (L312). The cause was
initialisation: each 7-day flow drove the hidden state to a saturated fixed
point, and an untrained policy's entropy fell from 1.07 to 0 within one
episode (L353). A closed-loop check (does the dose depend on the observation?)
is now a gate.
\item \textbf{The gate becomes wrong under missed visits.} Near-fixed schedules
scored \emph{higher} than closed-loop cells once observations were unreliable,
so excluding them biased the comparison (L632). We switched to
intention-to-treat and made MTD a reference arm to beat.
\item \textbf{Training chaos masquerades as architecture.} A $2\times10^{-7}$
float perturbation moved one ODE-RNN cell's mean $U$ from 132 to 256\,d.
Patient-only CIs ($[-20,-11]$) called the ODE-RNN worse than the LSTM, and
resampling seeds widened them to $[-43,+22]$ (L596).
\item \textbf{A mis-built shifted cohort.} Shifting $K$ after computing $y_0$
put 43--50\% of Tier-1 patients above capacity. Under full dose 14--18\% of
them diverged to NaN, and these were scored as censored, i.e.\ as successes
(L242). The progressor-only TTP median and TTP bought with extra drug misled
us in the same way, and both are why every TTP here comes with a censoring-aware
estimator and exposure.
\end{enumerate}

\section{Limitations}
One ODE family (LV), virtual patients only, and no claim about clinical
outcomes. The PQN settings are the published ones, untuned. Missed visits are
i.i.d., so the next gap is unknown when the policy doses. Scheduled gaps, where
continuous time has more to use, are untested. Mechanism shift and
within-episode drift (Tiers~2--3) are out of scope.

\section{Conclusion}
On this task, a memoryless burden threshold is a hard baseline to beat, and
continuous time did not beat it. When observations are unreliable, the
question is whether feedback pays at all before which architecture uses it
best. \pending{Final wording after S1.}

\begin{thebibliography}{9}
\bibitem{gallagher2024} K.~Gallagher, M.~A.~R. Strobl, D.~S. Park, et al.
Mathematical model-driven deep learning enables personalized adaptive therapy.
\emph{Cancer Research} 84(11):1929--1941, 2024. doi:10.1158/0008-5472.CAN-23-2040.
\bibitem{strobl2021} M.~A.~R. Strobl, J.~West, Y.~Viossat, M.~Damaghi,
M.~Robertson-Tessi, J.~S. Brown, R.~A. Gatenby, P.~K. Maini, A.~R.~A. Anderson. Turnover modulates the need for a cost of
resistance in adaptive therapy. \emph{Cancer Research} 81(4):1135--1147, 2021.
doi:10.1158/0008-5472.CAN-20-0806.
\bibitem{kiram2026} F.~M.~E. Kiram, I.~Youkana, R.~Saouli, G.~A. Susto,
L.~Kahloul. Recurrent deep reinforcement learning for chemotherapy control
under partial observability. arXiv:2605.02552, 2026.
\bibitem{rubanova2019} Y.~Rubanova, R.~T.~Q. Chen, D.~Duvenaud. Latent ODEs
for irregularly-sampled time series. \emph{NeurIPS}, 2019.
\bibitem{gallici2024} M.~Gallici, M.~Fellows, B.~Ellis, B.~Pou, I.~Masmitja,
J.~N. Foerster, M.~Martin. Simplifying deep temporal difference
learning. arXiv:2407.04811, 2024.
\end{thebibliography}
\end{document}
